\section{\texorpdfstring{The reciprocal turn $\golden^{-\lambda_{\log}}=\tfrac12$}{The reciprocal turn}}
\label{sec:spectral}

The transcendental period $\lambda_{\log}=\log_{\golden}2$ has two
directions.  Forward it closes the binary bridge
$\golden^{\lambda_{\log}}=2$ of clause~(iv) (\Cref{lem:binary}); read
backward it closes the reciprocal turn
$\golden^{-\lambda_{\log}}=\tfrac12$.  This value is not a thirteenth
clause but a point already inscribed in the identity, legible at three
of its nodes: it is the reciprocal of clause~(iv), the $p=-1$ node of
the universal identity~(v) (\Cref{thm:universal}), and the
multiplicative-line image of the conjugate root of the rotor
(clause~(i)).  We establish the three, then note where the value leads.

\begin{Lemma}[Reciprocal face of the binary bridge]\label{lem:inverse-turn}
$\golden^{-\lambda_{\log}}=\tfrac12$.
\end{Lemma}

\begin{proof}
By clause~(iv) (\Cref{lem:binary}), $\golden^{\lambda_{\log}}=2$; hence
$\golden^{-\lambda_{\log}}=2^{-1}=\tfrac12$.  In Lean this is the inverse
of \lean{mersenne\_bridge\_via\_lambda}.
\end{proof}

\begin{Corollary}[The reciprocal turn is the universal identity at $p=-1$]\label{cor:reciprocal-node}
The reciprocal turn $\golden^{-\lambda_{\log}}=\tfrac12$ is the value of
the universal identity~(v) (\Cref{thm:universal}) at $p=-1$.
\end{Corollary}

\begin{proof}
The universal identity reads $3\,\golden^{\sigma(p)}=2^{p}$ with
$\sigma(p)=p\,\lambda_{\log}-\log_{\golden}3$.  At $p=-1$ the binary side
is $2^{-1}=\tfrac12$, and the golden side reduces as
\[
  3\,\golden^{\sigma(-1)}
    =3\,\golden^{-\lambda_{\log}}\,\golden^{-\log_{\golden}3}
    =3\,\golden^{-\lambda_{\log}}\cdot\tfrac13
    =\golden^{-\lambda_{\log}},
\]
so~(v) at $p=-1$ reads $\golden^{-\lambda_{\log}}=\tfrac12$.
\end{proof}

The factor $3=M_{2}$ (clause~(vii)) that divides out here is the same
multiplier the identity carries at every~$p$---the cancellation
$3\,\golden^{-\log_{\golden}3}=1$ is the mechanism of~(v), not a
coincidence of $p=-1$; what is proper to $p=-1$ is only that the binary
side reads $\tfrac12$.  The reciprocal turn therefore sits inside the
central identity~(v), the identity that carries the ternary factor
$3=M_{2}$.

\begin{Remark}[Inversion image of the rotor]\label{rem:inversion-image}
Clause~(i) inscribes $\golden$ as a root and its inverse,
$\golden=\zeta_{10}+\zeta_{10}^{-1}$.  Clause~(iv) carries the same
pattern---a value and its inverse---to the multiplicative line: the pair
$(\zeta_{10},\zeta_{10}^{-1})$ on the unit circle corresponds to
$(\golden^{\lambda_{\log}},\golden^{-\lambda_{\log}})=(2,\tfrac12)$
on~$\mathbb{R}_{>0}$.  The rotor is closed under inversion
($e^{i\theta}\leftrightarrow e^{-i\theta}$); the reciprocal turn
$\tfrac12$ is, on~$\mathbb{R}_{>0}$, the analogue of the conjugate root
$\zeta_{10}^{-1}$, completing the inversion symmetry of the
multiplicative companion so that it mirrors the rotor it accompanies.
\end{Remark}

These three inscriptions are internal: each holds by the machinery of
the present identity, with no operator invoked.  This same value
$\tfrac12$ is where the companion operator paper anchors a self-adjoint
operator in the Hilbert--P\'olya programme (\Cref{ssec:euler-substrate});
that operator and its relation to the zeros of~$\zeta$ are developed
there.

